\documentclass[10pt,a4paper]{amsart}
\usepackage{amsmath,amssymb,amsthm,graphicx}
\usepackage{verbatim}
\newenvironment{reference}{\comment}{\endcomment}
%\newenvironment{reference}{}{}
\newenvironment{slogan}{\comment}{\endcomment}
\newenvironment{history}{\comment}{\endcomment}

% For commutative diagrams we use Xy-pic
\usepackage[all]{xy}

% We use 2cell for 2-commutative diagrams.
\xyoption{2cell}
% Theorem environments.
%
\theoremstyle{plain}
\newtheorem{theorem}[subsection]{Theorem}
\newtheorem{proposition}[subsection]{Proposition}
\newtheorem{lemma}[subsection]{Lemma}

\theoremstyle{definition}
\newtheorem{definition}[subsection]{Definition}
\newtheorem{example}[subsection]{Example}
\newtheorem{exercise}[subsection]{Exercise}
\newtheorem{situation}[subsection]{Situation}

\theoremstyle{remark}
\newtheorem{remark}[subsection]{Remark}
\newtheorem{remarks}[subsection]{Remarks}

\numberwithin{equation}{subsection}

% Macros
%
\def\lim{\mathop{\mathrm{lim}}\nolimits}
\def\colim{\mathop{\mathrm{colim}}\nolimits}
\def\Spec{\mathop{\mathrm{Spec}}}
\def\Hom{\mathop{\mathrm{Hom}}\nolimits}
\def\Ext{\mathop{\mathrm{Ext}}\nolimits}
\def\SheafHom{\mathop{\mathcal{H}\!\mathit{om}}\nolimits}
\def\SheafExt{\mathop{\mathcal{E}\!\mathit{xt}}\nolimits}
\def\Sch{\mathit{Sch}}
\def\Mor{\mathop{\mathrm{Mor}}\nolimits}
\def\Ob{\mathop{\mathrm{Ob}}\nolimits}
\def\Sh{\mathop{\mathit{Sh}}\nolimits}
\def\NL{\mathop{N\!L}\nolimits}
\def\CH{\mathop{\mathrm{CH}}\nolimits}
\def\proetale{{pro\text{-}\acute{e}tale}}
\def\etale{{\acute{e}tale}}
\def\QCoh{\mathit{QCoh}}
\def\Ker{\mathop{\mathrm{Ker}}}
\def\Im{\mathop{\mathrm{Im}}}
\def\Coker{\mathop{\mathrm{Coker}}}
\def\Coim{\mathop{\mathrm{Coim}}}

% Boxtimes
%
\DeclareMathSymbol{\boxtimes}{\mathbin}{AMSa}{"02}

%
% Macros for moduli stacks/spaces
%
\def\QCohstack{\mathcal{QC}\!\mathit{oh}}
\def\Cohstack{\mathcal{C}\!\mathit{oh}}
\def\Spacesstack{\mathcal{S}\!\mathit{paces}}
\def\Quotfunctor{\mathrm{Quot}}
\def\Hilbfunctor{\mathrm{Hilb}}
\def\Curvesstack{\mathcal{C}\!\mathit{urves}}
\def\Polarizedstack{\mathcal{P}\!\mathit{olarized}}
\def\Complexesstack{\mathcal{C}\!\mathit{omplexes}}
% \Pic is the operator that assigns to X its picard group, usage \Pic(X)
% \Picardstack_{X/B} denotes the Picard stack of X over B
% \Picardfunctor_{X/B} denotes the Picard functor of X over B
\def\Pic{\mathop{\mathrm{Pic}}\nolimits}
\def\Picardstack{\mathcal{P}\!\mathit{ic}}
\def\Picardfunctor{\mathrm{Pic}}
\def\Deformationcategory{\mathcal{D}\!\mathit{ef}}
\usepackage{cleveref}
\crefname{originaldefinition}{original definition}{original definitions}
\Crefname{originaldefinition}{Original definition}{Original definitions}
\makeatletter
\newlabel{divisors-definition-divisor-invertible-sheaf}{{0BE6}{source}{}{sourcepointer-1}{}}
\newlabel{divisors-definition-divisor-invertible-sheaf@cref}{{[originaldefinition][0][]0BE6}{[1][0][]source}{}{}{}}
\makeatother
\crefname{originallemma}{original lemma}{original lemmas}
\Crefname{originallemma}{Original lemma}{Original lemmas}
\makeatletter
\newlabel{divisors-lemma-Pic-projective-space-UFD}{{0BXJ}{source}{}{sourcepointer-2}{}}
\newlabel{divisors-lemma-Pic-projective-space-UFD@cref}{{[originallemma][0][]0BXJ}{[1][0][]source}{}{}{}}
\makeatother
\crefname{originallemma}{original lemma}{original lemmas}
\Crefname{originallemma}{Original lemma}{Original lemmas}
\makeatletter
\newlabel{divisors-lemma-normal-c1-injective}{{0BE8}{source}{}{sourcepointer-3}{}}
\newlabel{divisors-lemma-normal-c1-injective@cref}{{[originallemma][0][]0BE8}{[1][0][]source}{}{}{}}
\makeatother
\crefname{originallemma}{original lemma}{original lemmas}
\Crefname{originallemma}{Original lemma}{Original lemmas}
\makeatletter
\newlabel{divisors-lemma-open-subscheme-UFD}{{0BDA}{source}{}{sourcepointer-4}{}}
\newlabel{divisors-lemma-open-subscheme-UFD@cref}{{[originallemma][0][]0BDA}{[1][0][]source}{}{}{}}
\makeatother
\crefname{originalreference}{original reference}{original references}
\Crefname{originalreference}{Original reference}{Original references}
\makeatletter
\newlabel{lemma-blowup}{{0AGQ}{source}{}{sourcepointer-5}{}}
\newlabel{lemma-blowup@cref}{{[originalreference][0][]0AGQ}{[1][0][]source}{}{}{}}
\makeatother
\crefname{originalreference}{original reference}{original references}
\Crefname{originalreference}{Original reference}{Original references}
\makeatletter
\newlabel{lemma-blowup-regular}{{0AGR}{source}{}{sourcepointer-6}{}}
\newlabel{lemma-blowup-regular@cref}{{[originalreference][0][]0AGR}{[1][0][]source}{}{}{}}
\makeatother
\crefname{originallemma}{original lemma}{original lemmas}
\Crefname{originallemma}{Original lemma}{Original lemmas}
\makeatletter
\newlabel{more-algebra-lemma-regular-local-UFD}{{0AG0}{source}{}{sourcepointer-7}{}}
\newlabel{more-algebra-lemma-regular-local-UFD@cref}{{[originallemma][0][]0AG0}{[1][0][]source}{}{}{}}
\makeatother
\title{The Picard group of the blowup of the maximal ideal of a regular local surface}
\author{The Stacks Project Authors}
\date{stacks-project:resolve.tex:local-snapshot-20261012; GNU FDL 1.2 or later}
\begin{document}
\maketitle
\noindent\textit{Editorial note.} All delivered mathematical text is the author's own, copied from the pinned Stacks Project chapter file resolve.tex; no statement or proof step was rewritten, reordered or removed.\par\medskip
\noindent\textit{Editorial note.} The counted proof is the author's complete argument: the restriction of the twist to the exceptional divisor generating an infinite cyclic subgroup, the trivialization on the complement of the exceptional divisor, and the divisor computation identifying an arbitrary class.\par\medskip
\noindent\textit{Editorial note.} The label lemma-blowup-pic names the counted statement, and the proof environment immediately following it in the source is delivered verbatim as its proof.\par\medskip
\noindent\textit{Editorial note.} The counted statement is carried by the author's own lemma environment in the chapter source.\par\medskip
\noindent\textit{Editorial note.} Pre-open gate: python operations/agent-authoring/tools/precard\_gate.py --source resolve.tex --label lemma-blowup-pic --start 755 --end 783 --local-id 28104 returned RESULT=OK (L28104 free, occupancy FREE, public interval clear). Pre-assembly tag self-check: canonical tag resolved and not present in reference/public-items.jsonl.\par\medskip
\noindent\textit{Editorial note.} The statement and proof are copied from the pinned chapter file; the cited inputs are kept as external pointers with their canonical Stacks tags.\par\medskip
\noindent\textit{Allowed standard prerequisites (editorial).}\begin{itemize}
\item The Picard group of the projective line over a field and the restriction of the twist. Divisors, the lemma computing the Picard group of projective space over a unique factorization domain, together with the identification of the exceptional divisor and the twist on the blowup. Shows that the twist of the blowup generates an infinite cyclic subgroup of the Picard group through its restriction to the exceptional divisor.
\item Regular local rings are unique factorization domains and the punctured spectrum has trivial Picard group. More on Algebra, the lemma that a regular local ring is a unique factorization domain, and Divisors, the lemma on the Picard group of the complement of a closed point. Provides a trivialization of any invertible sheaf on the complement of the exceptional divisor.
\item Divisors of invertible sheaves and injectivity of the class group map on a normal scheme. Divisors, the definition of the divisor of a meromorphic section of an invertible sheaf, and the lemma that the class group map is injective for normal schemes. Turns the trivialization on the complement into a divisor supported on the exceptional divisor and identifies the sheaf with a twist.
\end{itemize}\medskip
\ifdefined\hypertarget\hypertarget{sourcepointer-1}{}\fi
\noindent\textit{Source reference.} Original definition 0BE6: sources/stacks-chow/tags.txt, source lines 14339--14339. Prerequisite pointer invoked at proof line(s) 779; it names a standard fact of a sibling chapter that the counted argument uses, and every step of the counted proof built on it is retained verbatim in the delivered body.\par\medskip
\ifdefined\hypertarget\hypertarget{sourcepointer-2}{}\fi
\noindent\textit{Source reference.} Original lemma 0BXJ: sources/stacks-chow/tags.txt, source lines 14982--14982. Prerequisite pointer invoked at proof line(s) 765; it names a standard fact of a sibling chapter that the counted argument uses, and every step of the counted proof built on it is retained verbatim in the delivered body.\par\medskip
\ifdefined\hypertarget\hypertarget{sourcepointer-3}{}\fi
\noindent\textit{Source reference.} Original lemma 0BE8: sources/stacks-chow/tags.txt, source lines 14341--14341. Prerequisite pointer invoked at proof line(s) 780; it names a standard fact of a sibling chapter that the counted argument uses, and every step of the counted proof built on it is retained verbatim in the delivered body.\par\medskip
\ifdefined\hypertarget\hypertarget{sourcepointer-4}{}\fi
\noindent\textit{Source reference.} Original lemma 0BDA: sources/stacks-chow/tags.txt, source lines 14308--14308. Prerequisite pointer invoked at proof line(s) 773; it names a standard fact of a sibling chapter that the counted argument uses, and every step of the counted proof built on it is retained verbatim in the delivered body.\par\medskip
\ifdefined\hypertarget\hypertarget{sourcepointer-5}{}\fi
\noindent\textit{Source reference.} Original reference 0AGQ: sources/stacks-resolve/pinned-original-source.tex, source lines 666--666. Prerequisite pointer invoked at proof line(s) 766; it names a standard fact of the same chapter that the counted argument uses, and every step of the counted proof built on it is retained verbatim in the delivered body.\par\medskip
\ifdefined\hypertarget\hypertarget{sourcepointer-6}{}\fi
\noindent\textit{Source reference.} Original reference 0AGR: sources/stacks-resolve/pinned-original-source.tex, source lines 727--727. Prerequisite pointer invoked at proof line(s) 781; it names a standard fact of the same chapter that the counted argument uses, and every step of the counted proof built on it is retained verbatim in the delivered body.\par\medskip
\ifdefined\hypertarget\hypertarget{sourcepointer-7}{}\fi
\noindent\textit{Source reference.} Original lemma 0AG0: sources/stacks-chow/tags.txt, source lines 13178--13178. Prerequisite pointer invoked at proof line(s) 770; it names a standard fact of a sibling chapter that the counted argument uses, and every step of the counted proof built on it is retained verbatim in the delivered body.\par\medskip
\begin{lemma}
\label{lemma-blowup-pic}
Let $(A, \mathfrak m, \kappa)$ be a regular local ring of dimension $2$.
Let $f : X \to S = \Spec(A)$ be the blowing up of $A$ in $\mathfrak m$.
Then $\Pic(X) = \mathbf{Z}$ generated by $\mathcal{O}_X(E)$.
\end{lemma}

\begin{proof}
Recall that $E = \mathbf{P}^1_\kappa$ has Picard group $\mathbf{Z}$
with generator $\mathcal{O}(1)$, see
Divisors, Lemma \ref{divisors-lemma-Pic-projective-space-UFD}.
By Lemma \ref{lemma-blowup} the invertible $\mathcal{O}_X$-module
$\mathcal{O}_X(E)$ restricts to $\mathcal{O}(-1)$. Hence
$\mathcal{O}_X(E)$ generates an infinite cyclic group in $\Pic(X)$.
Since $A$ is regular it is a UFD, see More on Algebra, 
Lemma \ref{more-algebra-lemma-regular-local-UFD}.
Then the punctured spectrum $U = S \setminus \{\mathfrak m\} = X \setminus E$
has trivial Picard group, see
Divisors, Lemma \ref{divisors-lemma-open-subscheme-UFD}.
Hence for every invertible $\mathcal{O}_X$-module $\mathcal{L}$
there is an isomorphism $s : \mathcal{O}_U \to \mathcal{L}|_U$.
Then $s$ is a regular meromorphic section of $\mathcal{L}$
and we see that $\text{div}_\mathcal{L}(s) = nE$ for some
$n \in \mathbf{Z}$
(Divisors, Definition \ref{divisors-definition-divisor-invertible-sheaf}).
By Divisors, Lemma \ref{divisors-lemma-normal-c1-injective}
(and the fact that $X$ is normal by Lemma \ref{lemma-blowup-regular})
we conclude that $\mathcal{L} = \mathcal{O}_X(nE)$.
\end{proof}

\end{document}
